\section{Introduction}

The OpenSees kernel was originally organized around object composition.
The finite-element model is represented by a \texttt{Domain} object,
while the procedures that number, assemble, and solve the governing
equations are delegated to a separate analysis composition.
Classes such as \texttt{AnalysisModel}, \texttt{ConstraintHandler},
\texttt{DOF\_Numberer}, \texttt{Integrator},
\texttt{EquiSolnAlgorithm}, and \texttt{LinearSOE} interact through
narrow abstract interfaces so that model objects remain decoupled from
numerical solution procedures.

The Tcl and Python front ends introduced a second architectural layer
whose responsibility is not finite-element analysis itself, but the
construction and coordination of analysis objects from a scripting
environment.
In the legacy implementation, this layer is realized through the OPS
interface, a collection of free functions operating on process-global
state.
That design departs from the composition discipline of the kernel in two
important ways.
First, runtime state is ambient rather than instance-scoped. 
Commands recover the active model, current load pattern, and active
analysis objects from global storage rather than from explicit arguments.
Second, the binding layer exposes analysis internals that were intended
to remain encapsulated within the analysis composition.
Functions such as \texttt{OPS\_GetAnalysisModel()},
\texttt{OPS\_GetAlgorithm()}, and \texttt{OPS\_GetSOE()} allow command
implementations to bypass the aggregation boundaries established in the
original architecture.

For interactive scripting this arrangement is serviceable.
For embedded use it is a fundamental limitation.
A process can safely host only one active OpenSees runtime, command
behavior depends on hidden program history, and independent software
components that rely on OpenSees cannot be isolated from one another.
These constraints obstruct shared-memory parallel evaluation, complicate
testing and reproducibility, and prevent the simulation kernel from
serving as a reentrant component within larger software systems.

This chapter presents a structured API that restores explicit ownership
and encapsulation at the language-binding boundary.
The central idea is to reify all binding-layer state as user-owned
runtime objects and to restrict command implementations to
capability-specific interfaces rather than unrestricted global accessors. 
Model-definition state, domain state, and analysis state are represented
by distinct objects with explicit lifetimes.
As a result, multiple OpenSees runtimes may coexist within one process,
the binding layer becomes reentrant, and the numerical kernel can be
embedded as a library component without altering the underlying
finite-element abstractions.

The contribution is therefore architectural rather than numerical.
The work does not introduce a new element formulation or a new solution
algorithm.
Its contribution is to restore, at the scripting boundary, the same
separation of concerns that the original kernel established within the
analysis core.


Several solutions have been proposed to address the architectural limitations in the OpenSeesPy interface \citep{marafi2022openseesapi,manousakis2022opensees,dong2022smartanalyze,talley2021opswrapper,cslotboom2021openseespytools,2022o3seespy}. 
However, all of these act as wrappers within the interpreter and do not address the underlying source of the problem.


\section{Legacy Binding Layer}

\subsection{Ambient Runtime State}

The legacy OPS layer stores the state of model construction and analysis
in process-global variables distributed across the interpreter source.
These variables include the active domain, the current model builder,
the active load pattern and time series, and the objects that define the
analysis stack.
Command implementations do not receive these dependencies explicitly.
Instead, they query ambient state through free functions in the OPS
interface.

This design has two immediate consequences.
The first is hidden sequential coupling.
The meaning of a command depends not only on its arguments, but also on
which commands have previously executed.
For example, the interpretation of a \texttt{load} command depends on
which pattern object was most recently installed as the current pattern.
Likewise, analysis commands act on whichever domain and analysis objects
happen to occupy the global slots at the time of the call.
The second consequence is that state ownership becomes implicit.
Nothing in a command signature reveals which model it mutates or which
analysis objects it depends upon.

\subsection{Violation of Encapsulation}

The original OpenSees architecture places objects such as
\texttt{AnalysisModel}, \texttt{Integrator},
\texttt{EquiSolnAlgorithm}, and \texttt{LinearSOE} inside an analysis
composition with well-defined responsibilities.
The OPS layer weakens these boundaries by publishing accessors that
return raw pointers to such objects.
A command implementation can therefore reach past the public interface of
the analysis object and manipulate internal collaborators directly.

This capability leak defeats the very decoupling that the kernel
architecture was designed to provide.
An extension that conceptually requires only model-building services
acquires access to solver internals.
An analysis command can depend on objects that are not visible in its
signature and not governed by a local ownership relation.
The result is an interface that is easy to extend in an ad hoc manner
but difficult to reason about compositionally.

\subsection{Single-Instance Semantics}

Because the OPS layer stores the active runtime in process-global state,
the interpreter effectively supports only one live model and one live
analysis stack per process.
Constructing a new model replaces the previously active one.
Reconfiguring the solver stack mutates shared state rather than the state
of a particular analysis object.
This precludes safe coexistence of independent OpenSees-based libraries
within one address space and forces ensemble workflows toward process
isolation even when shared-memory concurrency would otherwise be natural.

\section{Structured API}

\subsection{Runtime Object Model}

The redesigned interface replaces ambient interpreter state with explicit
runtime objects.
The state previously hidden in the OPS layer is partitioned into three
categories.

The first category is model-definition state.
This includes registries of materials, sections, coordinate
transformations, time series, and other reusable objects referenced by
tag during model construction.

The second category is domain state.
This is the evolving finite-element model itself, consisting of nodes,
elements, constraints, load patterns, and committed response history.

The third category is analysis state.
This includes the objects that define and execute a particular solution
procedure, such as the constraint handler, degree-of-freedom numberer,
integrator, solution algorithm, convergence test, and system of
equations.

Each category is represented by an explicit object with a well-defined
owner and lifetime.
No process-global variables are used to recover these objects.

\subsection{Capability-Restricted Commands}

In the new interface, a command operates only on the object that is
explicitly passed to it.
A model-building routine receives access to model-definition and domain
services, but not to the internal collaborators of an active analysis.
An analysis-construction routine receives the specific handles required
to assemble an analysis object, but not unrestricted access to all other
runtime state.
This restriction preserves encapsulation by ensuring that commands depend only on the
capabilities they actually require.

The practical effect is that the binding layer no longer bypasses the
aggregation structure of the kernel.
Internal analysis objects remain internal.
The only way to alter an analysis is through the public interface of the
analysis runtime that owns them.

\subsection{Consequences}

Once runtime state is instance-scoped, multiple OpenSees runtimes may
coexist in a single process without interpreter-induced interference.
The binding layer becomes reentrant because all mutable state is held by
explicit objects rather than ambient global storage.
This permits library-grade embedding, simplifies test isolation, and
allows concurrent evaluation of disjoint models.

The same restructuring also improves compatibility with sensitivity and
differentiation workflows.
Because model state and analysis state are explicit program values, host
software can construct, retain, duplicate, and schedule them without
relying on hidden process history.
The resulting interface is therefore suitable not only for scripting, but
also for integration into optimization, reliability, and monitoring
systems that require explicit control over object lifetimes and
execution order.